\textbf{Phần 1. Trắc nghiệm lựa chọn}
\setcounterpageref{ex}{0}

\begin{ex} %[2D7N36]
Số gồm $9$ chục và $4$ đơn vị viết là
\choice
{$49$}
{$904$}
{$90$}
{\True $94$}
\loigiai{
Số đó là $94$.\\
\textbf{Đáp án của câu hỏi là D}
}
\end{ex}

\begin{ex} %[2D7N36]
Số liền sau của số $99$ là
\choice
{$98$}
{\True $100$}
{$101$}
{$90$}
\loigiai{
$99 + 1 = 100$.\\
\textbf{Đáp án của câu hỏi là B}
}
\end{ex}

\begin{ex} %[2D7H36]
Kết quả của phép tính $56 + 37$ là
\choice
{$92$}
{$94$}
{\True $93$}
{$83$}
\loigiai{
$56 + 37 = 93$.\\
\textbf{Đáp án của câu hỏi là C}
}
\end{ex}

\begin{ex} %[2D7H36]
Kết quả của phép tính $91 - 47$ là
\choice
{\True $44$}
{$43$}
{$45$}
{$54$}
\loigiai{
$91 - 47 = 44$.\\
\textbf{Đáp án của câu hỏi là A}
}
\end{ex}

\begin{ex} %[2D7H36]
Đường gấp khúc gồm 3 đoạn thẳng có độ dài lần lượt là $8\text{ cm}, 10\text{ cm}, 12\text{ cm}$. Độ dài đường gấp khúc là
\choice
{$28\text{ cm}$}
{\True $30\text{ cm}$}
{$32\text{ cm}$}
{$34\text{ cm}$}
\loigiai{
$8 + 10 + 12 = 30\text{ cm}$.\\
\textbf{Đáp án của câu hỏi là B}
}
\end{ex}

\begin{ex} %[2D7H36]
Một can đựng $15\text{ l}$ dầu, một can khác đựng $20\text{ l}$ dầu. Cả hai can đựng số lít dầu là
\choice
{$30\text{ l}$}
{$40\text{ l}$}
{$45\text{ l}$}
{\True $35\text{ l}$}
\loigiai{
$15 + 20 = 35\text{ l}$.\\
\textbf{Đáp án của câu hỏi là D}
}
\end{ex}

\begin{ex} %[2D7V36]
Lớp 2A trồng được 38 cây, lớp 2B trồng được nhiều hơn lớp 2A là 15 cây. Lớp 2B trồng được số cây là
\choice
{$52$ cây}
{$54$ cây}
{\True $53$ cây}
{$55$ cây}
\loigiai{
$38 + 15 = 53$ cây.\\
\textbf{Đáp án của câu hỏi là C}
}
\end{ex}

\begin{ex} %[2D7V36]
Tìm $x$ biết: $x + 28 = 74$.
\choice
{\True $46$}
{$45$}
{$47$}
{$48$}
\loigiai{
$x = 74 - 28 = 46$.\\
\textbf{Đáp án của câu hỏi là A}
}
\end{ex}

\begin{ex} %[2D7V36]
Tính giá trị biểu thức: $75 - 28 + 19$.
\choice
{$64$}
{$68$}
{$70$}
{\True $66$}
\loigiai{
$47 + 19 = 66$.\\
\textbf{Đáp án của câu hỏi là D}
}
\end{ex}

\begin{ex} %[2D7V36]
Nếu hôm nay là thứ Năm ngày 12 thì thứ Năm tuần trước là ngày
\choice
{\True $5$}
{$4$}
{$6$}
{$7$}
\loigiai{
$12 - 7 = 5$.\\
\textbf{Đáp án của câu hỏi là A}
}
\end{ex}

\begin{ex} %[2D7V36]
Tính nhanh: $45 + 18 + 15 + 22$.
\choice
{$98$}
{\True $100$}
{$102$}
{$90$}
\loigiai{
$(45 + 15) + (18 + 22) = 60 + 40 = 100$.\\
\textbf{Đáp án của câu hỏi là B}
}
\end{ex}

\begin{ex} %[2D7C36]
Tìm số có hai chữ số có tổng các chữ số bằng 14, biết chữ số hàng chục hơn chữ số hàng đơn vị là 2.
\choice
{$77$}
{$95$}
{\True $86$}
{$68$}
\loigiai{
Chữ số hàng chục: $(14 + 2) : 2 = 8$.\\
Chữ số hàng đơn vị: $8 - 2 = 6$. Số đó là $86$.\\
\textbf{Đáp án của câu hỏi là C}
}
\end{ex}

\textbf{Phần 2. Trắc nghiệm đúng sai}
\setcounterpageref{ex}{0}

\begin{ex} %[2D7H36]
Xét tính đúng sai về kiến thức học kì 1:
\choiceTF
{\True Số lớn nhất có hai chữ số là 99.}
{\True $48 + 36 = 84$.}
{Hình tứ giác có 3 cạnh.}
{\True Hình tứ giác có 4 cạnh.}
\loigiai{
\begin{itemize}
    \item Số lớn nhất có 2 chữ số là 99: Đúng.
    \item $48 + 36 = 84$: Đúng.
    \item Tứ giác có 4 cạnh (không phải 3). Ý 3 Sai, Ý 4 Đúng.
\end{itemize}
\begin{itemize}
    \item a) \textbf{Đúng}.
    \item b) \textbf{Đúng}.
    \item c) \textbf{Sai}.
    \item d) \textbf{Đúng}.
\end{itemize}
}
\end{ex}

\begin{ex} %[2D7V36]
Một trang trại có 52 con bò và 38 con trâu. Xét tính đúng sai:
\choiceTF
{\True Trang trại có tất cả 90 con trâu và bò ($52 + 38 = 90$).}
{\True Số bò nhiều hơn số trâu là 14 con ($52 - 38 = 14$).}
{Nếu bán đi 12 con bò thì số bò còn lại ít hơn số trâu.}
{\True Khi đó còn $52 - 12 = 40$ con bò (vẫn nhiều hơn 38 con trâu).}
\loigiai{
\begin{itemize}
    \item Tổng: $52 + 38 = 90$: Đúng.
    \item Bò hơn trâu: $52 - 38 = 14$: Đúng.
    \item Bò còn 40 con $> 38$ con trâu. Ý 3 Sai, Ý 4 Đúng.
\end{itemize}
\begin{itemize}
    \item a) \textbf{Đúng}.
    \item b) \textbf{Đúng}.
    \item c) \textbf{Sai}.
    \item d) \textbf{Đúng}.
\end{itemize}
}
\end{ex}

\textbf{Phần 3. Trả lời ngắn}
\setcounterpageref{ex}{0}

\begin{ex} %[2D7H36]
Tính tổng: $49 + 38$.
\shortans{87}
\loigiai{
\textbf{Đáp số:} $87$\\
$49 + 38 = 87$.
}
\end{ex}

\begin{ex} %[2D7H36]
Tính hiệu: $85 - 39$.
\shortans{46}
\loigiai{
\textbf{Đáp số:} $46$\\
$85 - 39 = 46$.
}
\end{ex}

\begin{ex} %[2D7V36]
Tính độ dài đường gấp khúc gồm 3 đoạn thẳng có độ dài là $7\text{ cm}, 13\text{ cm}, 15\text{ cm}$ theo xăng-ti-mét.
\shortans{35}
\loigiai{
\textbf{Đáp số:} $35$\\
$7 + 13 + 15 = 35\text{ cm}$.
}
\end{ex}

\begin{ex} %[2D7V36]
Tìm $x$ biết: $x - 36 = 47$.
\shortans{83}
\loigiai{
\textbf{Đáp số:} $83$\\
$x = 47 + 36 = 83$.
}
\end{ex}

\begin{ex} %[2D7V36]
Một cửa hàng có $75\text{ kg}$ đường, đã bán đi $48\text{ kg}$. Hỏi cửa hàng còn lại bao nhiêu ki-lô-gam đường?
\shortans{27}
\loigiai{
\textbf{Đáp số:} $27$\\
$75 - 48 = 27\text{ kg}$.
}
\end{ex}

\begin{ex} %[2D7C36]
Tính nhanh tổng: $16 + 27 + 34 + 23$.
\shortans{100}
\loigiai{
\textbf{Đáp số:} $100$\\
$(16 + 34) + (27 + 23) = 50 + 50 = 100$.
}
\end{ex}
